\subsection{道成分の排除と閉路上の係数の向き}
\label{pc12:section}

高分母 support の次数は，その core における高分母 rays の本数に等しい．
従って，一高分母の負係数と定理~\ref{td12:one-negative} を組み合わせると，
異なる高分母を許した場合にも，次数二以下の成分を制限できる．

\begin{theorem}[高分母の道成分の不存在]\label{pc12:no-path}
現在の全平面五次の構成条件の下で，大域作用適合 atlas が存在するとする．
高分母 support
\[
 H=\{ij:q_{ij}>1\}\subset K_5
\]
は，辺を持つ path を一つの連結成分として持たない．
分母の一致や上限を仮定せず，node 数の非負性と全 node budget も用いない．
\end{theorem}

\begin{proof}
path 成分の頂点を順に $i_0,i_1,\ldots,i_m$，$m\geq1$ とする．
各葉は一高分母 core なので，命題~\ref{qh9:isolated-local} により
\[
 D^{i_0}_{i_1}<0,\qquad D^{i_m}_{i_{m-1}}<0.
\]
対向端相殺から $D^{i_1}_{i_0}>0$ となる．
$m=1$ なら既に矛盾する．$m\geq2$ のとき，$i_1$ は二高分母 core であるから，
定理~\ref{td12:one-negative} により，もう一つの高分母係数は
$D^{i_1}_{i_2}<0$ でなければならない．
同じ議論を順に繰り返すと，最後の辺で $D^{i_m}_{i_{m-1}}>0$ を得る．
これは最後の葉の負係数に反する．
\end{proof}

\begin{proposition}[閉路に沿う非零係数]\label{pc12:cycle-orientation}
同じ条件の下で，$H$ の一つの連結成分が cycle であるとする．
その全ての高分母係数は非零であり，各 core には正と負の係数が一つずつある．
各辺を負の係数を持つ端から正の係数を持つ端へ向けると，
全辺の向きは一つの有向閉路となる．
特に，$H$ の最大次数が二以下で $H$ が辺を持つなら，
辺を持つ成分は一つだけで，それは四角形または五角形である．
\end{proposition}

\begin{proof}
cycle の頂点数と辺数を $m$ とする．各 core は二高分母を持つので，
定理~\ref{td12:one-negative} により，負の高分母係数が少なくとも一つある．
従って負の係数の総数は少なくとも $m$ である．
一方，各辺では対向端相殺により，両係数が零であるか，
一方が負でもう一方が正である．従って負の係数は各辺に高々一つであり，
総数は高々 $m$ である．等号が成立するため，各 core に負の係数はちょうど一つで，
零の係数を持つ辺はない．各頂点の入次数と出次数は共に一であるから，向きは閉路に沿って揃う．

最大次数二以下の有限単純 graph の，辺を持つ連結成分は path または cycle である．
path は定理~\ref{pc12:no-path}，三角形は定理~\ref{ls11:no-triangle-component} で排除される．
$H\subset K_5$ なので残るのは四角形または五角形であり，二つのそのような成分は共存できない．
\end{proof}

\begin{proposition}[木成分の中心への予備的な還元]\label{pc12:tree-reduction}
同じ条件の下で，$H$ の辺を持つ連結成分が tree であるなら，
その graph は $K_{1,3}$，または $K_{1,3}$ の一辺を一度細分した五頂点の tree に限られる．
いずれの場合にも，唯一の三価 core における三つの高分母係数は全て正でなければならない．
\end{proposition}

\begin{proof}
五頂点以下の tree で最大次数二以下のものは path なので排除済みである．
最大次数四のものは $K_{1,4}$ に限られるが，各葉の負係数と対向端相殺により，
中心の四つの係数は全て正となる．これは row sums から従う
$\sum_jD^i_j=0$ に反する．
最大次数三のものは，記載した二種類に限られる．
葉から中心へ進む各道において，定理~\ref{pc12:no-path} の証明と同じ符号伝播を用いると，
中心に到達する全ての高分母係数が正になる．
\end{proof}

これらは連結成分についての主張である．別の高分母辺が接続する path 部分 graph や，
三高分母以上を持つ core の任意の係数配置を排除してはいない．
後の定理~\ref{fo13:no-tree} では，葉からの係数の大きさを用いて全ての木成分を排除する．
一般三高分母の全高係数の正性については，独立した未解決の局所問題として残る．
